\documentclass[12pt,a4paper]{article}

\usepackage[german]{babel}
\usepackage[utf8]{inputenc}
\usepackage[T1]{fontenc}
\usepackage{enumerate}
\usepackage[absolute]{textpos}
\usepackage[a4paper, top=2cm, left=2cm, bottom=2cm, right=2cm]{geometry}
\usepackage{parskip}

\usepackage{graphicx}

\begin{document}
\thispagestyle{empty}

\begin{textblock*}{\paperwidth}(1.6cm,0.48cm) %{width}(horizontal, vertical)
\includegraphics[scale=0.67]{Logo_Mathematik_CMYK.pdf}
\end{textblock*}

\vspace*{1.5cm}

\begin{center}
\Huge{\textbf{Informations on the examination process}}\\
\large{Mathematical Olympiad 2020}
\end{center}

\bigskip

The following guidelines prevail for any examination organised by the Mathematical Olympiad.

\section{Permitted material}
\begin{itemize}
\item Permitted material includes pens and geometry tools (compass, ruler, protractor).
\item The use of calculators and formularies is forbidden. The use of personal sheets of paper is not allowed during the exam.
\item The consumption of drinks and snacks is allowed during the exam. In considerations of other contestants one should forgo food which consumption is noisy or diffuses a strong smell.
\end{itemize}

\section{Proceedings of the exam}
\begin{itemize}

\item Solutions and sketches must be written on the official paper provided by the invigilators. It is always possible during the exam to ask for more paper.

\item Questions regarding the problem statements are permitted at any time. All questions must be written down on a separate sheet of paper and handed to an invigilator.

\item Visits to the toilet during the exam is allowed only with the authorization  of an invigilator.

\end{itemize}

\section{Writing and marking}
\begin{itemize}
\item You must not write on the back of the sheets. Two different problems must never be treated on the same sheet of paper.
\item The contestant's full name, the number of the exercise, the page number and the total number of pages used for that exercise must be written down on every sheet of paper, e.g. "Sir Andrew Wiles, G1, p.2/4".
\item Every step of the solution must be justified. Every invoked mathematical result must be properly referenced.
\item After reception of the corrected exam it is possible, within the announced dead-line, to ask for clarifications regarding the correction. If justified, the correction may be re-evaluated.
\end{itemize}
\end{document} 